% Generated draft: requires mathematical and source review.
\documentclass[11pt]{article}
\usepackage[margin=25mm]{geometry}
\usepackage{fontspec}
\setmainfont{Latin Modern Roman}
\usepackage{amsmath,amssymb,mathtools,unicode-math}
\setmathfont{Latin Modern Math}
\usepackage{enumitem,needspace,xurl,hyperref}
\setlength{\parindent}{0pt}
\setlength{\parskip}{5pt}
\setlength{\emergencystretch}{3em}
\newcommand{\N}{\mathbb{N}}
\newcommand{\Z}{\mathbb{Z}}
\newcommand{\Q}{\mathbb{Q}}
\newcommand{\R}{\mathbb{R}}
\newcommand{\C}{\mathbb{C}}
\newcommand{\F}{\mathbb{F}}
\newcommand{\A}{\mathbb{A}}
\newcommand{\PP}{\mathbb{P}}
\newcommand{\E}{\mathbb{E}}
\newcommand{\eps}{\varepsilon}
\newcommand{\dd}{\,\mathrm{d}}
\begin{document}
\section*{UnsolvedMath ID: 20002364}
\textbf{Source problem code:} catalogue entry
\subsection{Definitions and mathematical statement}
Let $K$ be an algebraically closed field and $F=K(Z)$, with $Z$ transcendental over $K$. A nontrivial valuation ring of a field $M$ means a proper subring $\mathcal O\subsetneq M$ containing $1$ such that for every $x\in M^\times$, either $x\in\mathcal O$ or $x^{-1}\in\mathcal O$. Shlapentokh's question asks whether there exist such $K$ and such a valuation ring in $F$, or in a finite extension of $F$, which is definable. In the pure ring-language reading with parameters this asks, for example, for a formula $\varphi(x,\bar a)$ with $\bar a\in F$ defining a proper valuation ring of $F$. The original wording does not specify the characteristic, parameter allowance, quantifier complexity or whether the valuation must be trivial on $K$. Its phrase ``or a finite extension'' also leaves the ambient definability convention unstated.

The supplied aid gives the following parameter obstruction for valuations trivial on the constants. If $v$ is a nontrivial valuation on $F$ whose restriction to $K^\times$ is zero, then its valuation ring is not definable in the pure ring language using only finitely many parameters from $K$. Indeed the places of $K(Z)/K$ correspond to the points of $\PP^1(K)$, and the fractional-linear action of $\operatorname{PGL}_2(K)$, by automorphisms fixing $K$ pointwise, moves every such valuation ring. A definable set with parameters from $K$ would be invariant under all these automorphisms. Thus a definition of a nontrivial $K$-trivial valuation ring must use a nonconstant parameter. This obstruction is not asserted for valuations whose restrictions to $K$ are nontrivial.

The second reported reduction is independent of that constant-triviality restriction. Let $L/F$ be finite and $\mathcal O_w$ a nontrivial valuation ring of $L$ definable in its pure field structure, allowing parameters. Then
\[
 \mathcal O_v=\mathcal O_w\cap F
\]
is a nontrivial valuation ring of $F$ definable with parameters from $F$. To formulate the transfer, choose an $F$-basis $e_1=1,e_2,\ldots,e_n$ of $L$, and name in $F$ its finitely many multiplication structure constants and the coordinates of the parameters in the definition of $\mathcal O_w$. Interpret $L$ on $F^n$ using these constants; substitute $(x,0,\ldots,0)$ for the argument of the defining formula. This gives the contraction formula in $F$. If the definition in $L$ is existential, so is its translated definition in $F$, since each quantified $L$-variable is replaced by finitely many $F$-variables. Nontriviality is preserved: a valuation trivial on $F$ has only the trivial extension to an algebraic extension. The descent need not preserve being parameter-free or definable from constants alone, because the basis constants and parameter coordinates may be nonconstant.

Accordingly the finite-extension option does not avoid the parameter-allowing base-field existence problem under this interpretation. The aid supplies these reductions, not an example solving that problem. Its statement that a positive definition must name a nonconstant element is justified by the symmetry argument only in the $K$-trivial case.
\subsection{Short English statement}
Seek a definable nontrivial valuation ring in a rational function field over some algebraically closed field, possibly after a finite extension. The aid excludes definitions from constants alone for valuations trivial on the constants and shows that a finite-extension example descends to the base field when parameters are allowed.
\subsection{Sources}
\begin{itemize}
\item Shlapentokh, Question 27, \emph{Problems related to ``Extensions of Hilbert's Tenth Problem''}, American Institute of Mathematics; moderators B. Poonen and T. Scanlon, notes by J. Demeyer, as identified in the supplied verification. \url{https://aimath.org/WWN/hilberts10th/hilberts10th.pdf}.
\item Local UnsolvedMath registry, permanent ID 20002364, source code \texttt{catalogue entry}, canonical location \texttt{aim-logic-notes.json}, \texttt{notes[139]}, record dated 14 August 2026. The original literature field is empty.
\item Supplied research attempt 1: reduction, problem status at run open, rigor labelled proved, novelty confidence low. The symmetry and finite-extension descent claims are machine-generated research aid, not peer reviewed, requiring independent verification before citation; no separate publication for these reductions is supplied.
\end{itemize}
\end{document}
